\documentclass[12pt, ukrainian]{article}
\setlength\textwidth{190mm} \setlength\hoffset{-25mm}
\setlength\textheight{277mm} \setlength\voffset{-28mm}
%\documentclass[10pt, ukrainian]{article}
%\setlength\textwidth{190mm} \setlength\textheight{277mm}
%\setlength\voffset{-38mm} \setlength\hoffset{-38mm}
%\setlength\parindent{0mm}
\pagestyle{empty}
\usepackage[T2A]{fontenc}
\usepackage[utf8]{inputenc}
\usepackage[ukrainian]{babel}
\usepackage{graphicx}
\usepackage{caption}
\usepackage{multicol}
\usepackage{verbatim}
\DeclareCaptionFormat{nocaption}{}
\captionsetup{format=nocaption,aboveskip=0pt,belowskip=0pt}
\usepackage[Export]{adjustbox}
\adjustboxset{max size={0.9\linewidth}{0.9\paperheight}}
\begin{document}
{{\begin {small}\par
{ \center  Київський національний університет імені Тараса Шевченка \par}
{\parbox {5 cm}{Освітня програма \hfill}{\parbox {7 cm}{Системний аналіз \hfill}}\parbox {6 cm}{\hfill Семестр 2}}\par
{\parbox {5 cm} {Навчальний предмет \hfill}\parbox {7 cm}{Програмування \hfill}\parbox {6cm}{\hfill}\par}\vspace{-7pt} \end{small}}{\center Екзаменаційний білет \No \textbf
{ 25 }\par}}
{
\begin {large}
УВАГА!

У письмовій роботі мають зазначатися всі відповіді, а також проміжні міркування, що передували їх отриманню.

Проміжні міркування можуть бути якоюсь арифметикою, відслідковуванням зна\-чень змінних для деяких завдань,
короткими поясненнями, допоміжною інформацією, що виникає під час обходу графів, робочею чернеткою,
тобто тим, що відображає хід розв’язання або пояснює отриману відповідь.
Це є ОБОВ’ЯЗКОВИМ для кожного запитання.

Без наявності доречних проміжних міркувань відповіді зараховуватись як правиль\-ні не будуть.
Писати багато та красиво оформлено не треба. Коротко, але по темі. Можли\-во, навіть чернетка.

Якщо на аркуші буде чернетка разом з відповідями --- це ОК.

\newpage

{\begin {center}\textbf{Завдання 1}\end {center}}\par
Що буде виведено за виконання?\par
\begin{center}
	\adjustimage{max size={0.9\linewidth}{0.9\paperheight}}
	{../pict/25_1.png}
\end{center}
\newpage

{\begin {center}\textbf{Завдання 2}\end {center}}\par
Що буде виведено за виконання?\par
 (повідомлення інтерпретатора про помилку позначати error)\par
\begin{center}
	\adjustimage{max size={0.9\linewidth}{0.9\paperheight}}
	{../pict/25_2.png}
\end{center}
\newpage

{\begin {center}\textbf{Завдання 3}\end {center}}\par
Що буде виведено за виконання?\par
 (повідомлення інтерпретатора про помилку позначати error)\par
\begin{center}
	\adjustimage{max size={0.9\linewidth}{0.9\paperheight}}
	{../pict/25_3.png}
\end{center}
\newpage

{\begin {center}\textbf{Завдання 4}\end {center}}\par
Що буде виведено за виконання?\par
 (повідомлення інтерпретатора про помилку позначати error)\par
\begin{center}
	\adjustimage{max size={0.9\linewidth}{0.9\paperheight}}
	{../pict/25_4.png}
\end{center}
\newpage

{\begin {center}\textbf{Завдання 5}\end {center}}\par
Що буде виведено за виконання?\par
 (повідомлення інтерпретатора про помилку позначати error)\par
\begin{center}
	\adjustimage{max size={0.9\linewidth}{0.9\paperheight}}
	{../pict/25_5.png}
\end{center}
\newpage

{\begin {center}\textbf{Завдання 6}\end {center}}\par
Що буде виведено за виконання?\par
 (повідомлення інтерпретатора про помилку позначати error)\par
\begin{center}
	\adjustimage{max size={0.9\linewidth}{0.9\paperheight}}
	{../pict/25_6.png}
\end{center}
\newpage

{\begin {center}\textbf{Завдання 7}\end {center}}\par
Що буде виведено за виконання?\par
 (повідомлення інтерпретатора про помилку позначати error)\par
\begin{center}
	\adjustimage{max size={0.9\linewidth}{0.9\paperheight}}
	{../pict/25_7.png}
\end{center}
\newpage

{\begin {center}\textbf{Завдання 8}\end {center}}\par
Що буде виведено за виконання?\par
 (повідомлення інтерпретатора про помилку позначати error)
%Оригинал был утерян. Требует отладки!\par
\begin{center}
	\adjustimage{max size={0.9\linewidth}{0.9\paperheight}}
	{../pict/25_8.png}
\end{center}
\newpage

{\begin {center}\textbf{Завдання 9}\end {center}}\par
Що буде виведено за виконання?\par
 (повідомлення інтерпретатора про помилку позначати error)\par
\begin{center}
	\adjustimage{max size={0.9\linewidth}{0.9\paperheight}}
	{../pict/25_9.png}
\end{center}
\newpage

{\begin {center}\textbf{Завдання 10}\end {center}}\par
Для графа на рисунку з вершини 4 запущено алгоритм пошуку в ширину, що будує кістякове дерево. Списки суміжності вершин переглядаються алгоритмом за зростанням міток вершин.\par
\begin{center}
	\adjustimage{max size={0.9\linewidth}{0.9\paperheight}}
	{../Graphs/Traversals/349.png}
\end{center}
\par Які з наведених ребер входять до складу отриманого кістякового дерева?\par
1) (0, 6)\par
2) (1, 2)\par
3) (1, 5)\par
4) (2, 5)\par
5) (0, 1)\par
6) (2, 3)\par
{\vskip 15pt} Номери відповідних ребер записати за зростанням через пробіл.\par
\newpage


\end {large}
}
{\smallskip \begin {small} Затверджено на засіданні кафедри теоретичної кібернетики, протокол \No 6 від   29.01.2022 р.\par\smallskip\smallskip
{\parbox {6.5 cm}{Зав. кафедри \dotfill Ю.В.~Крак}}\hspace {0.7cm}{\hbox to 6.5 cm{ Екзаменатор \dotfill Т.О.~Карнаух}} \end{small}}
\end{document}
